\section{Methods}
\label{sec:methods}

% Present the solution and the EXPECTATIONS you have from it (course writing recs).

\subsection{Data}
% CommonsenseQA. State explicitly: evaluation uses the VALIDATION split (~1,221
% items) because test-set answer keys are withheld for the leaderboard.

\subsection{Answer-side translation}
% How answer choices are translated; backend (MT vs LLM — see CLAUDE.md decisions);
% how meaning preservation is checked (optional human validation on a subset).

\subsection{Conditions}
\label{sec:conditions}
% Define the Q-A language conditions and the two extensions.
% \begin{table}[t]
%   \centering
%   \begin{tabular}{lll}
%     \toprule
%     Condition & Question & Answers \\
%     \midrule
%     EN-EN & English & English (baseline) \\
%     EN-X  & English & X \\
%     X-X   & X       & X \\
%     X-EN  & X       & English \\
%     \bottomrule
%   \end{tabular}
%   \caption{Q--A language conditions. X is a target language (Hebrew, Arabic,
%            Spanish/French). Extensions: mixed answer sets; partial question-noun
%            translation.}
%   \label{tab:conditions}
% \end{table}

\subsection{Models and prompting}
% Evaluated model(s); zero/few-shot; the prompt template. State the answer-selection
% convention: the model emits a CHOICE LETTER (A--E), never the answer text, so that
% scoring is language-invariant and flips are not a string-matching artifact.

\subsection{Metrics}
% Accuracy per condition; flip rate between condition pairs (the core diagnostic);
% statistical significance (e.g. McNemar for paired accuracy on the same items).

\subsection{Reproducibility}
% temperature=0, fixed seed, cached raw outputs, per-run manifest (model id+snapshot,
% prompt-template hash, variant hash, decoding params). Point to src/ + configs/.
